\chapter{KADABRA}
\section{Overview and Problem Formalization}

While exact algorithms like Brandes' compute Betweenness Centrality (BC) for all nodes in $O(mn)$ time, this worst-case complexity becomes prohibitive for large-scale real-world networks. KADABRA shifts the paradigm from exact computation to sampling-based approximation.

\subsection{The Stochastic Paradigm}
The core premise relies on a probabilistic interpretation of betweenness centrality.
 If we select a source node $s$ and a target node $t$ uniformly at random,
  and subsequently choose a single shortest path $\pi$ between them uniformly at random from all available $st$-shortest paths,
   the betweenness centrality $bc(v)$ of a node $v$ is exactly equal to the probability that $\pi$ passes through $v$:
\begin{equation}
    \Pr(v \in \pi) = bc(v)
\end{equation}

KADABRA leverages this by sampling a sequence of $\tau$ random shortest paths $\pi_1, \dots, \pi_\tau$ according to this exact distribution.
 For each sampled path, an empirical indicator variable $X_i(v)$ is recorded:
\begin{equation}
    X_i(v) = \begin{cases} 1 & \text{if } v \in \pi_i \text{ and } v \neq s, t \\ 0 & \text{otherwise} \end{cases}
\end{equation}

The algorithm's empirical estimate of betweenness centrality, denoted as $\tilde{b}(v)$, is simply the sample mean of these indicators:
\begin{equation}
    \tilde{b}(v) := \frac{1}{\tau} \sum_{i=1}^{\tau} X_i(v)
\end{equation}

By definition, this estimator is unbiased, yielding an expected value $\mathbb{E}[\tilde{b}(v)] = bc(v)$.

\subsection{Objective and Error Formalization}
The objective of KADABRA is to minimize the required number of sampled paths $\tau$ while providing strict statistical guarantees on the quality of the approximation across all vertices. The user specifies two global tolerance parameters:
\begin{itemize}
    \item An \textbf{absolute error bound} $\lambda > 0$.
    \item A \textbf{confidence parameter} $\delta \in (0, 1)$.
\end{itemize}

The standard problem task is to ensure that, upon termination, the empirical estimates for all nodes fall within a $1-\delta$ simultaneous confidence interval. Formally, the algorithm guarantees that:
\begin{equation}
    \Pr(\forall v \in V, |\tilde{b}(v) - bc(v)| \le \lambda) \ge 1 - \delta
\end{equation}

\subsection{The Adaptive Sampling Challenge}
A na\"ive approximation algorithm fixes a static sample size $\tau$ in advance using worst-case bounds derived from standard concentration inequalities (such as Hoeffding's or Chernoff's bounds).
 However, this uniformly over-samples nodes with low centrality, which require significantly fewer paths to approximate accurately.

KADABRA optimizes this by utilizing \textbf{adaptive sampling}, meaning the total number of samples $\tau$ is treated as a dynamic random variable (a stopping time) determined on the fly based on the specific paths observed during execution.

The technical challenge—and a major point of departure from prior literature—is that evaluating a
 stopping condition based on the data introduces a subtle \textbf{stochastic dependence} between the termination time $\tau$ 
 and the correctness of the estimator. KADABRA resolves this by using martingale theory to ensure that checking bounds dynamically 
 does not compromise the mathematical strictness of the $(\lambda, \delta)$-guarantee.


 \section{Algorithmic Innovations in KADABRA}

KADABRA achieves its massive speedup over previous state-of-the-art approximation algorithms by optimizing both the time required to generate an individual path sample and the total number of samples required to guarantee precision.

\subsection{Balanced Bidirectional BFS (bb-BFS)}
Most existing betweenness approximation algorithms generate a sample path by picking a random source node $s$ and executing a standard single-source Breadth-First Search (BFS) to explore the graph. In the worst-case, this demands $\Theta(|E|)$ time per sample. 

KADABRA introduces a \textbf{balanced bidirectional BFS} (bb-BFS) routine to optimize path generation. Given a random source $s$ and target $t$, the algorithm simultaneously grows a forward BFS tree from $s$ and a backward BFS tree from $t$. To ensure the search remains computationally balanced, the algorithm keeps track of the boundary volumes:
\begin{itemize}
    \item Let $\Gamma^{l_s}(s)$ denote the set of vertices at distance $l_s$ from $s$.
    \item Let $\Gamma^{l_t}(t)$ denote the set of vertices at distance $l_t$ from $t$.
\end{itemize}

At each step, the algorithm chooses to expand from the tree with the smaller boundary volume, minimizing the immediate edge-traversal cost:
\begin{equation}
\text{Expand from } s \iff \sum_{v \in \Gamma^{l_s}(s)} \text{deg}(v) \le \sum_{u \in \Gamma^{l_t}(t)} \text{deg}(u)
\end{equation}

The bidirectional search terminates as soon as the boundaries of the two trees ``touch each other'' (i.e., a common edge or node is discovered).
 A valid $st$-shortest path is then constructed by connecting a uniformly chosen random path from $s$ to the meeting boundary,
  and tracking back from the boundary to $t$. 

Crucially, the authors prove that on realistic random complex networks (such as Configuration Models or Rank-1 Inhomogeneous Random Graphs),
 bb-BFS cuts down path generation time from worst-case linear time $\Theta(|E|)$ to $|E|^{\frac{1}{2}+o(1)}$ with high probability.

\subsection{Rigorous Adaptive Stopping Condition}
To minimize the number of samples $\tau$, KADABRA evaluates a stopping condition dynamically on the fly.
 Rather than tracking error from a predefined global count,
  KADABRA sets up an absolute maximum sample limit $\omega$ (derived from worst-case VC-dimension bounds) and computes deterministic target
   allocation confidences $(\delta_L^{(v)}, \delta_U^{(v)})$ for each node $v$ during an initial brief profiling phase.

At each step $\tau$, the algorithm evaluates lower and upper error margins ($f$ and $g$) based on the empirical centrality computed up to that point.
 The mathematical bounds are derived as:
\begin{equation}
f(\tilde{b}(v), \delta_L^{(v)}, \omega, \tau) = \frac{1}{\tau}\ln\frac{1}{\delta_L^{(v)}}\left(\frac{1}{3} - \frac{\omega}{\tau} + \sqrt{\left(\frac{1}{3} - \frac{\omega}{\tau}\right)^2 + \frac{2\tilde{b}(v)\omega}{\ln(1/\delta_L^{(v)})}}\right)
\end{equation}
\begin{equation}
g(\tilde{b}(v), \delta_U^{(v)}, \omega, \tau) = \frac{1}{\tau}\ln\frac{1}{\delta_U^{(v)}}\left(\frac{1}{3} + \frac{\omega}{\tau} + \sqrt{\left(\frac{1}{3} + \frac{\omega}{\tau}\right)^2 + \frac{2\tilde{b}(v)\omega}{\ln(1/\delta_U^{(v)})}}\right)
\end{equation}

The sampling loop dynamically terminates as soon as the loop condition maps to \texttt{true},
 checking if $f \le \lambda$ and $g \le \lambda$ for all relevant vertices. By framing the sequence of sample path indicators as a martingale,
  the algorithm ensures that evaluating this stopping rule dynamically over a stochastic running time $\tau$ does not break the target $(1-\delta)$ probability guarantee.

\section{Theoretical Guarantees: The Contrast with GNNs}

The foundational justification for employing KADABRA over other modern approximation methodologies lies in its rigorous mathematical guarantees. When evaluating a sampling-based approach against learning-based models like Graph Neural Networks (GNNs), this section highlights the trade-offs between provable statistical bounds and inductive heuristics.

\subsection{Rigorous Martingale-Based Correctness}
The core mathematical achievement of KADABRA is proving that its dynamic stopping condition preserves strict probabilistic safety. In standard adaptive sampling, terminating an algorithm based on the observed data introduces a stochastic dependence that can bias the final estimates. 

By modeling the sequence of path indicator variables as a martingale $Z^\tau$ and applying McDiarmid's martingale concentration inequality, the authors prove that KADABRA delivers an absolute $(\lambda, \delta)$-approximation upon termination:
\begin{equation}
\Pr\Big(\forall v \in V, \, |\tilde{b}(v) - bc(v)| \le \lambda\Big) \ge 1 - \delta
\end{equation}

Whether the algorithm stops early ($\tau < \omega$) via its adaptive triggers, or reaches the worst-case sample bound ($\omega$), the mathematical certainty of the confidence interval remains completely unbroken.

% \subsection{Core Architectural Differences from GNN Approaches}
% This formal framework sets up a stark contrast with GNN-based betweenness centrality estimators, highlighting a fundamental divergence in design paradigms:

% \begin{itemize}
%     \item \textbf{Data-Agnostic vs. Data-Dependent Performance:} KADABRA requires absolutely no training data or offline optimization phase. Its error bounds hold strictly for \textit{any} topology, whereas GNNs rely on the assumption that the target graph shares structural distributions with the training dataset.
%     \item \textbf{Absolute Guarantees vs. Heuristic Approximations:} KADABRA yields a provable upper bound on the estimation error controlled explicitly by the user via $\lambda$ and $\delta$. GNNs minimize an empirical loss function and output predictions that lack strict, per-node error bounds on unseen graphs.
%     \item \textbf{Execution Burden (Runtime vs. Training Time):} KADABRA shifts all computational weight to the execution phase by actively traversing the graph using balanced bidirectional BFS. Conversely, GNNs transfer this heavy computational burden to an offline training phase; once trained, the GNN can predict centrality scores for an entire graph in a single forward pass ($O(|V| + |E|)$), bypassing the need for explicit path exploration entirely.
% \end{itemize}

\subsection{Adaptive Top-$k$ Centrality Selection}
A significant limitation of previous betweenness approximation frameworks is their architectural rigidity: they require computing a uniform $\lambda$-approximation across \textit{all} nodes to construct a reliable top-$k$ ranking. This introduces massive computational overhead, as identical precision is wasted on low-centrality nodes that are mathematically far from the top-$k$ boundary. KADABRA resolves this by decoupling individual node confidence intervals, allowing larger margins of error for nodes whose centrality values are structurally well-separated.

Given a user-defined threshold $k$ and absolute tolerance $\lambda$, KADABRA guarantees that upon termination, for every vertex $v$, the algorithm either:
\begin{enumerate}
    \item Establishes the exact correct positional rank of $v$.
    \item Formally guarantees that $v \notin \text{Top-}k$.
    \item Provides an empirical estimate such that $|bc(v) - \tilde{b}(v)| \le \lambda$.
\end{enumerate}

To implement this on the fly, KADABRA replaces the global, uniform checking condition with a dynamic boundary validation filter. At each sampling iteration $\tau$, the vertices are sorted in descending order of their empirical centrality estimates to form a tentative sequence $v_1, v_2, \dots, v_n$. The algorithm is allowed to safely terminate early if and only if the following strict separation criteria are simultaneously met:

\begin{itemize}
    \item \textbf{Internal Top-$k$ Boundary Separation:} For every internal index $i \in [1, k]$, if the bounds $f(v_i)$ or $g(v_i)$ still exceed $\lambda$, the algorithm checks if their empirical confidence intervals overlap with their immediate neighbors:
    \begin{equation}
    \tilde{b}(v_{i-1}) - f(v_{i-1}) \ge \tilde{b}(v_i) + g(v_i) \quad \text{and} \quad \tilde{b}(v_i) - f(v_i) \ge \tilde{b}(v_{i+1}) + g(v_{i+1})
    \end{equation}
    If these lower and upper bounds cross, the true positions are not yet statistically distinct, forcing the algorithm to sample more paths.
    
    \item \textbf{External Boundary Exclusion:} For every node outside the top ranking $i \in [k+1, n]$, if the bounds $f(v_i)$ or $g(v_i)$ exceed $\lambda$, the algorithm ensures that the node cannot breach the Top-$k$ barrier:
    \begin{equation}
    \tilde{b}(v_k) - f(v_k) \ge \tilde{b}(v_i) + g(v_i)
    \end{equation}
    This strictly confirms that $v_i$ cannot have a true betweenness centrality higher than the $k$-th most central node.
\end{itemize}

This adaptive relaxation allows KADABRA to skip massive volumes of unnecessary shortest-path sampling. In practice, it enables the first successful approximation of top-$k$ centralities on graphs featuring hundreds of millions of edges (such as massive citation or collaboration networks) in less than an hour.